%%%%%%%%%%  scan idx 43  |  folio 32  %%%%%%%%%%

\begin{center}
\underline{Expos\'e VI}

\vspace{1em}
\underline{\underline{TEORÍA FORMAL DE LAS DEFORMACIONES}}

\vspace{1em}
\underline{por D. S. Rim}
\end{center}

1. \underline{Teorema de existencia formal}

Recordemos la noción de categoría cofibrada (S.G.A. 1 IV). Sea
$\underline C$ una categoría fija. Una categoría cofibrada $\underline A$
sobre $\underline C$ es, por definición, una categoría $\underline A$
junto con un funtor $P:\underline A\longrightarrow\underline C$ tal que:

\underline{Ax. 1.} Para todo morfismo $f$ de $\underline C$ y todo objeto
$a$ de $\underline A$ tal que $P(a)$ sea la fuente de $f$, existe un
$f$-morfismo cocartesiano $\alpha:a\longrightarrow b$. Así, para todo
$f$-morfismo $\alpha':a\longrightarrow b'$, existe un único
$T$-morfismo $\tau:b\longrightarrow b'$ en $\underline A$ tal que
$\alpha'=\tau\alpha$, donde $T$ es el blanco de $f$.

\underline{Ax. 2.} Un compuesto de morfismos cocartesianos en
$\underline A$ es de nuevo cocartesiano.

Las propiedades siguientes de una categoría cofibrada
$P:\underline A\longrightarrow\underline C$ se deducen inmediatamente de
las definiciones.

(Cancelación) Sean
$a\overset{\alpha}{\underset{\beta}{\rightrightarrows}}a'$ dos morfismos
cocartesianos tales que $P(\alpha)=P(\beta)$. Si existe un morfismo
cocartesiano $b\xrightarrow{\gamma}a$ tal que
$\alpha\gamma=\beta\gamma$, entonces $\alpha=\beta$.

(Factorización) Dados morfismos cocartesianos
$a\xrightarrow{\alpha}a'$ y $a\xrightarrow{\beta}a''$, existe un
morfismo cocartesiano $a'\xrightarrow{\gamma}a''$ tal que
$\gamma\alpha=\beta$ si y solo si existe un morfismo
$f:P(a')\longrightarrow P(a'')$ en $\underline C$ tal que
$fP(\alpha)=P(\beta)$. Además, $\gamma$ es único.

Sea $\underline A$ una categoría cofibrada sobre $\underline C$. Para cada
objeto $S$ de $\underline C$, denotemos por $\underline A(S)$ la
subcategoría de $\underline A$ cuyos objetos son los objetos $a$ para los
cuales $P(a)=S$, y en la que
\[
 \Hom_{\underline A(S)}(a,a')=
 \{\alpha\in\Hom_{\underline A}(a,a')\mid P(\alpha)=\id_S\}.
\]
Un grupoide cofibrado sobre $\underline C$ es una categoría cofibrada
$\underline A$ sobre $\underline C$ tal que $\underline A(S)$

%%%%%%%%%%  scan idx 44  |  folio 33  %%%%%%%%%%

es un grupoide para todo objeto $S$ de $\underline C$. Recordemos que un
grupoide es una categoría en la que todo morfismo es un isomorfismo. Así,
un grupoide cofibrado sobre $\underline C$ es una categoría cofibrada
$P:\underline A\longrightarrow\underline C$ en la que todo morfismo de
$\underline A$ es cocartesiano. Además, para todo
$\alpha\in\operatorname{Fl}(\underline A)$, $\alpha$ es un isomorfismo
en $\underline A$ si y solo si $P(\alpha)$ es un isomorfismo.

En lo que resta de esta sección, fijemos un anillo local noetheriano
completo $\Lambda$ con cuerpo residual $k$, y sea $\underline C_\Lambda$
la categoría de las $\Lambda$-álgebras locales artinianas con cuerpo
residual $k$, cuyos morfismos son homomorfismos locales sobre $\Lambda$.
A partir de 1.2, se supondrá siempre que un grupoide cofibrado
$\underline A$ sobre $\underline C_\Lambda$ satisface
$\underline A(k)=\{e\}$, donde $\{e\}$ denota la categoría en la que
$\Hom_{\{e\}}(a,b)$ consta de un único elemento para cada par de objetos
$a,b$ de $\{e\}$. Un morfismo $a'\xrightarrow{\alpha}a$ en
$\underline A$ se llamará simplemente $\underline C_\Lambda$-morfismo si
$P(a')=P(a)$ y $P(\alpha)=1_{P(a)}$. Por abuso de lenguaje, un morfismo
$a'\xrightarrow{\alpha}a$ en $\underline A$ se llamará sobreyectivo si
$P(\alpha):P(a')\longrightarrow P(a)$ es sobreyectivo en
$\underline C_\Lambda$.

\underline{Ej. 1.1(a).} Sea
$F:\underline C\longrightarrow(\text{categorías})$ un funtor.
Definamos la categoría $\underline F$ como sigue. Un objeto de
$\underline F$ es un par $(S,a)$ con
$S\in\mathrm{ob}(\underline C)$ y $a\in\mathrm{ob}F(S)$; un morfismo
$(S,a)\longrightarrow(S',a')$ es un par $(f,u')$, donde
$f:S\longrightarrow S'$ es un morfismo en $\underline C$ y
$u':F(f)a\longrightarrow a'$ es un morfismo en $F(S')$. El compuesto de
\[
 (S,a)\xrightarrow{(f,u')}(S',a')
 \xrightarrow{(g,u'')}(S'',a'')
\]
es
\[
 (gf,\,u''\mathbin{\cdot}F(g)u'):(S,a)\longrightarrow(S'',a'').
\]
Respecto del funtor
$P:\underline F\longrightarrow\underline C$ dado por
$(S,a)\longmapsto S$, $\underline F$ se convierte en una categoría
cofibrada sobre $\underline C$. Si $F(S)$ es un grupoide para todo
$S\in\mathrm{ob}(\underline C)$, entonces $\underline F$ es un grupoide
cofibrado sobre $\underline C$. En particular, un funtor
$F:\underline C\longrightarrow(\text{grupos})$ da lugar a un grupo
cofibrado, y un funtor $F:\underline C\longrightarrow(\mathrm{Ens})$
da lugar a un grupoide discreto cofibrado. Una transformación natural
$\varphi:F\longrightarrow G$, donde $F,G:\underline C\longrightarrow
(\text{categorías})$, induce un funtor cocartesiano
$\varphi:\underline F\longrightarrow\underline G$ sobre $\underline C$.

\underline{Ej. 1.1(b).} Sean $X$ e $Y$ esquemas sobre $\Lambda$
provistos de un

%%%%%%%%%%  scan idx 45  |  folio 34  %%%%%%%%%%

isomorfismo fijo
\[
 X\otimes_\Lambda k\xleftrightarrow{\ \tau\ }Y\otimes_\Lambda k.
\]
Entonces el funtor
$\Isom_{X,Y}:\underline C_\Lambda\longrightarrow(\text{grupoides})$
dado por
\[
 \Isom_{X,Y}(S)=
 \left\{X\otimes_\Lambda S\xleftrightarrow{\ \sigma\ }
 Y\otimes_\Lambda S\;\middle|\;
 \sigma\otimes_\Lambda k=\tau\right\}
\]
da lugar a un grupoide cofibrado
$P:\underline{\Isom}_{X,Y}\longrightarrow\underline C_\Lambda$.

\underline{Ej. 1.1(c).} Sea $X_0$ un esquema fijo sobre $k$, y sea
$\underline M_{X_0}$ la categoría de las deformaciones de $X_0$ sobre
$\underline C_\Lambda$ (véase \S4). Así, un objeto de
$\underline M_{X_0}$ es un par $(R,X)$, donde
$R\in\underline C_\Lambda$ y $X$ es una deformación de $X_0$ a $R$.
Entonces el funtor
$P:\underline M_{X_0}\longrightarrow\underline C_\Lambda$ dado por
$(R,X)\mathrel{\mapstochar\rightsquigarrow} R$ define un grupoide cofibrado sobre
$\underline C_\Lambda$.

\underline{Ej. 1.1(d).} Para todo grupoide $X$, denotemos por $[X]$ el
grupoide discreto constituido por las clases de isomorfía de objetos de
$X$. Si $P:\underline F\longrightarrow\underline C_\Lambda$ es un
grupoide cofibrado, obtenemos un funtor
$[\underline F]:\underline C_\Lambda\longrightarrow(\mathrm{Ens})$
dado por $[\underline F](S)=[\underline F(S)]$, y por tanto un grupoide
discreto cofibrado que, por abuso de notación, se denota de nuevo por
$[\underline F]$.
